% Originally: content/week-04.tex, \subsection{The mean value theorem} --- promoted to a section

% Quelle: Diego Torres Tejeda/Notes - 09.03 - Diego Torres Tejeda.pdf, p. 8
\section{The mean value theorem}
\label{sec:mean_value_theorem}

\begin{ainote}
The mean value theorem is used twice in this chapter --- in the official hint to
\Cref{ex:4.6} and in its solution --- but Corsin never states it. Diego Torres Tejeda's notes
state it, together with the warning that follows, which is what \cref{ex:4.6} actually turns on.
\end{ainote}

\begin{theorem}[Mean value theorem]
\label{thm:mean_value_theorem}
Let $f : U \to \mathbb{R}$ be differentiable, $x \in U$ and $h \in \mathbb{R}^n$ with
$x+th \in U$ for all $t \in [0,1]$. Then there is $t_0 \in (0,1)$ such that, writing
$\xi := x+t_0h$,
\[f(x+h)-f(x) = \partial_hf(\xi).\]
\end{theorem}

\begin{proof}
Apply the one-variable mean value theorem to $\varphi(t) := f(x+th)$ on $[0,1]$, which is
differentiable with $\varphi'(t) = \partial_hf(x+th)$ by the chain rule along a path
(\cref{sec:chain_rule_along_path}). It gives $t_0 \in (0,1)$ with
$\varphi(1)-\varphi(0) = \varphi'(t_0)$, which is the claim.
\end{proof}

\begin{importantremark}[It fails for vector-valued $f$]
\label{rem:mvt_fails_vector_valued}
The hypothesis \qt{$f : U \to \mathbb{R}$} is essential --- \textbf{the theorem is false the
moment the target is $\mathbb{R}^m$ with $m \geq 2$}, and this is exactly what
\cref{ex:4.5} below probes.
% Corrected: this pointed at \cref{ex:4.6}, which applies the *scalar* theorem to $u(x+t\nu)$
% and never leaves $\mathbb{R}$. The problem that probes the vector-valued failure is 4.5 ---
% which was not in the built document when the sentence was written.

Applying the scalar result to each component $f_1,\dots,f_m$ separately does produce points
$\xi_1,\dots,\xi_m$ with
\[f(x+h)-f(x) = \begin{pmatrix}\partial_hf_1(\xi_1) \\ \vdots \\ \partial_hf_m(\xi_m)\end{pmatrix},\]
but the $\xi_i$ are \emph{different points} in general, and no single $\xi$ need work for all
components at once. The standard witness is the \emph{circle}
$\gamma(t) := (\cos t, \sin t)$ on $[0,2\pi]$: here $\gamma(2\pi)-\gamma(0) = 0$, while
$\gamma'(t) = (-\sin t,\cos t)$ has $|\gamma'(t)| = 1$ and so never vanishes --- there is no
$\xi$ at all. Geometrically this is entirely reasonable: a walker who returns to the starting
point must have had zero \emph{average} velocity, but need never have stood still. In one
dimension \qt{returning} forces a turning point; in two it does not.

What survives for vector-valued $f$ is the \emph{mean value inequality},
$|f(x+h)-f(x)| \leq \sup_{t\in[0,1]}\lVert Df_{x+th}\rVert\,|h|$, which is what one actually
uses in practice.
\end{importantremark}

\begin{center}
% Generator: Claude Opus 5 (High) --- FIG-W04-06: the scalar theorem and its vector-valued
% failure, side by side.
% Left: f(x) = 0.35 + 1.5x - 0.55x^2 on [0.15, 2.4]. Chord slope between a=0.15 and b=2.4 is
% f'(xi) with f'(x) = 1.5 - 1.1x, and (f(b)-f(a))/(b-a) = 1.5 - 1.1*(a+b)/2 = 1.5 - 1.4025
% = 0.0975, so xi = (1.5 - 0.0975)/1.1 = 1.275 -- the midpoint (a+b)/2, as it must be for a
% parabola. Marked there.
% Right: the circle t -> (cos t, sin t); start = end, yet the velocity has length 1 throughout.
\begin{tikzpicture}[>=Latex, scale=1.0]

  % ---------- Left: the scalar mean value theorem ----------
  \begin{scope}
    \draw[->] (-0.2,0) -- (3.0,0) node[right, font=\footnotesize] {$x$};
    \draw[->] (0,-0.2) -- (0,2.5) node[above, font=\footnotesize] {$y$};

    \draw[blue!75!black, thick, domain=0.15:2.55, samples=70, smooth]
      plot (\x, {0.35 + 1.5*\x - 0.55*\x*\x});

    % Chord from a = 0.35 (y = 0.8076) to b = 2.4 (y = 0.7820); slope -0.0125.
    % For a parabola the mean value point is the midpoint, xi = (0.35+2.4)/2 = 1.375,
    % where f = 1.3727 and f' = 1.5 - 1.1*1.375 = -0.0125 -- equal to the chord slope.
    \draw[red!80!black, thick] (0.35,0.8076) -- (2.4,0.7820);
    \fill (0.35,0.8076) circle (1.2pt) node[below, font=\scriptsize] {$x$};
    \fill (2.4,0.7820) circle (1.2pt) node[below right, font=\scriptsize] {$x+h$};

    \draw[green!50!black, thick] (0.70,1.3811) -- (2.05,1.3643);
    \fill[green!50!black] (1.375,1.3727) circle (1.4pt);
    \node[green!50!black, font=\scriptsize, above] at (1.375,1.44) {$\xi$};

    \node[font=\scriptsize, align=center, anchor=north] at (1.4,-0.42)
      {\textbf{scalar $f$:} some $\xi$ has\\ tangent parallel to the chord};
  \end{scope}

  % ---------- Right: the vector-valued failure ----------
  \begin{scope}[shift={(5.6,1.15)}]
    \draw[->] (-1.5,0) -- (1.55,0) node[above right, font=\footnotesize] {$x$};
    \draw[->] (0,-1.5) -- (0,1.6) node[above, font=\footnotesize] {$y$};

    \draw[blue!75!black, thick] (0,0) circle (1);
    \fill[red!80!black] (1,0) circle (1.6pt);
    % Label placed below the axis: the velocity arrow at t = 0 points straight up from (1,0),
    % and the x-axis label sits just above the axis, so both sides are occupied.
    \draw[red!80!black, thin] (1,0) -- (1.18,-0.55);
    \node[red!80!black, font=\scriptsize, anchor=north west] at (1.10,-0.52)
      {$\gamma(0)=\gamma(2\pi)$};

    % velocity vectors, always of length 1, never zero
    \foreach \a in {0,72,144,216,288}{
      \draw[green!50!black, very thick, ->]
        ({cos(\a)},{sin(\a)}) -- ++({-0.52*sin(\a)},{0.52*cos(\a)});
    }
    \node[green!50!black, font=\scriptsize, anchor=east] at (-1.1,0.95)
      {$|\gamma'|=1$};

    \node[font=\scriptsize, align=center, anchor=north] at (0,-1.57)
      {\textbf{vector-valued $\gamma$:} $\gamma(2\pi)-\gamma(0)=0$,\\ yet $\gamma'$ never
       vanishes --- no $\xi$ exists};
  \end{scope}
\end{tikzpicture}
\end{center}

% Source: Corsin Nick/Class Notes/Week 4.pdf, p. 7
\begin{exercise}[Differential of $A \mapsto A\transp A$]
\label{ex:transpose_diff}
We identify the inner product space of real-valued matrices $\mathbb{R}^{n\times n}$ with inner
product
\[\langle A, B\rangle = \sum_{i,j=1}^{n} A_{ij}B_{ij} = \Tr(A\transp B)\]
for $A, B \in \mathbb{R}^{n\times n}$, with the euclidean space
$\mathbb{R}^{n^2}\cong\mathbb{R}^{n\times n}$, with the standard inner product. Then, consider
the function
\[F : \mathbb{R}^{n\times n}\to\mathbb{R}^{n\times n}, \qquad A \mapsto A\transp A.\]
Compute the differential at the identity matrix $\id \in \mathbb{R}^{n\times n}$ applied to
$X \in \mathbb{R}^{n\times n}$:
\[DF_{\id}(X) = \partial_X F(\id).\]
\end{exercise}

% Kept inline: solved directly in Corsin Nick/Class Notes/Week 4.pdf, p. 7.
\begin{exercisesolution}[\cref{ex:transpose_diff}]
We use $(*)$:
\[
\begin{aligned}
DF_{\id}(X) &= \left.\frac{d}{ds}\right|_{s=0} F(\id+sX) \\
&= \left.\frac{d}{ds}\,(\id+sX)\transp(\id+sX)\right|_{s=0} \\
&= \left.\frac{d}{ds}\left[\id + s(X\transp+X) + s^2X\transp X\right]\right|_{s=0} \\
&= X\transp + X
\end{aligned}
\]
\end{exercisesolution}


% --- exercise relocated here from the dissolved sheet 4 ---

% Originally: content/exercise-sheets/week-04.tex, problem 4.5 --- never previously built
% Source: exercises/Ex4_Analysis2_eng.pdf

\begin{exercise}[4.5 --- Mean value for vector-valued functions \textnormal{(optional)}]
\label{ex:4.5}
Let $f \in C^1(\mathbb{R},\mathbb{R}^m)$ for $m > 1$. Is it true that there is $t \in [0,1]$ such
that
$$f(1) - f(0) = Df_t(1) = \begin{pmatrix} f_1'(t) \\ \vdots \\ f_m'(t) \end{pmatrix}?$$
Prove it or provide a counterexample.

\textit{Official hint:} try $f(t) = (\sin(2\pi t), \cos(2\pi t))$.
\end{exercise}
